\section{Introduction}
\label{sec:intro}

As large reasoning models have come to solve hard problems through long solutions, distillation that teaches small models with those solutions has become widely used \citep{guo2025deepseekr1,muennighoff2025s1}. Small students, however, perform worse when they learn the long solutions of reasoning models \citep{li2025small,luo2025valley}, so recent work uses an interpolated teacher, which mixes a reasoning model and a short-answer instruct model with ratio $\lambda$, to adjust solution length and choose the point that suits the student \citep{lan2026midistillation}.

Such an interpolated teacher stops thinking altogether at some $\lambda$ \citep{wu2026interpolation,lan2025spectrum}, but what then remains in the student's weights is unknown.

Interpolation studies have examined only the outputs of the interpolated teacher, and one study reports that no such switch occurs \citep{yang2026system}. Distillation work selects $\lambda$ by the student's accuracy and, on the student side, has shown only that long solutions train the student along more concentrated directions than short ones \citep{lan2026midistillation}. Other studies read the weights or internals of distilled models, but they compare a single distilled model with its original model rather than comparing students while varying the teacher \citep{ward2025rank1,hu2026transcoder,venhoff2026base}. Moreover, in long solutions the way of thinking and the length move together, so outputs alone cannot tell which of the two has changed.

We therefore look at the weights that training leaves behind instead of the outputs, and ask two questions: what changes in the student's weights when the teacher stops thinking, and does the changed part carry the way of thinking?

Because full fine-tuning spreads its changes widely and is hard to read, we confine training to a LoRA adapter and use it as a measuring instrument. In the nine combinations of three teacher pairs and three students, we compare adapters trained while varying only the teacher through $\lambda$. We take the overlap between adapters that differ only in their training data as the baseline, so that only the changes caused by the teacher are isolated. Whether a model thinks is measured, for the teacher, by the think rate (the share of responses that use the \texttt{</think>} tag) and, for the student, by the frequency of reflective markers in its solutions. We then remove, graft or amplify the region we find and examine the outputs and the internal signals read with a sparse autoencoder (SAE). Our contributions are as follows.

\begin{itemize}
  \item We show in all three pairs that, at exactly the $\lambda$ where the teacher stops thinking, the student's adapter also sharply changes the directions along which it updates the weights (\S\ref{sec:switch}). Each pair follows its own teacher to a different $\lambda$, and the alignment holds whether the teacher is measured by the think rate or by reflective markers. Before that point, the adapters always use in common a shared region made up of 5--8 of the 16 directions.
  \item We show that the shared region acts like a switch that turns thinking on and off, but that grafting it alone does not carry the solution to the end (\S\ref{sec:carries}). Removing it reduces reflective markers to 0.2--10\% of their original level and makes answers 4--5 times shorter, whereas removing the same amount elsewhere in the adapter does not turn thinking off as much. Unlike a length-matched cut of the original model, removal turns thinking off and also gives higher accuracy. Amplifying the region with the norm held fixed does not increase thinking, and grafting it into a short-answer student makes it start thinking but not finish the solution. Internal signals read with the SAE mostly lead to the same conclusion.
  \item We show that the shared region also appears outside the training problems and in distillation without an interpolated teacher (\S\ref{sec:general}). Removing the region turns thinking off on GSM8K and GPQA-Diamond, problems outside the MATH data used for training, and a substantial part of the region is recovered in reasoning-distillation adapters trained without an interpolated teacher.
\end{itemize}
